\section{Results: diffuse scattering and stacking disorder in PbI\texorpdfstring{$_2$}{2}}
\label{sec:results_pbi2}

Figure~\ref{fig:pbi2_diffuse_cylinders} shows a CVD-grown PbI$_2$ film and a
crop spanning the $m=0$ and $m=1$ Bragg rods.  The $m=1$ rod contains
measurable intensity between the integer Bragg maxima, whereas the $m=0$ rod
is dominated by sharp peaks and finite-thickness oscillations.  This
rod-dependent redistribution of intensity is the experimental motivation for
the stacking-disorder model.  In a two-dimensional powder, random in-plane
rotation maps each in-plane reflection family onto a reciprocal-space cylinder
at fixed $Q_R$.  Perfect stacking concentrates intensity at discrete $Q_z$
values on a cylinder; loss of stacking coherence redistributes part of that
intensity along the same cylinder.

\begin{figure}[H]
  \centering
  \begin{minipage}{0.78\linewidth}
    \centering
    \includegraphics[width=\linewidth]{figures/results_pbi2/PbI2_detector.png}
  \end{minipage}
  \caption{Measured PbI$_2$ diffraction image (right) and an enlarged region
  of interest (left).  The direct-beam position defines the vertical
  $m=0$ trajectory.  Diffuse intensity is visible between the ordered maxima
  on the neighboring $m=1$ rod, while the $m=0$ profile remains insensitive
  to the lateral registry shifts that generate stacking contrast on other
  rods.}
  \label{fig:pbi2_diffuse_cylinders}
\end{figure}

\subsection{Direction-resolved transition-matrix model}
\label{sec:pbi2_transition_model}

The structural unit is one I--Pb--I trilayer.  Its state is written
$\alpha=(r,\sigma)$, where $r\in\{0,1,2\}$ denotes one of three basal
registries and $\sigma\in\{+,-\}$ denotes the two trilayer orientations.  The
six states are therefore
\begin{equation}
  \mathcal{S}=\left[0F_+,1F_+,2F_+,0F_-,1F_-,2F_-\right].
  \label{eq:pbi2_six_states}
\end{equation}
For a fixed $m$ rod, a forward registry step
$\mathbf{s}=(1/3,2/3)$ contributes
\begin{equation}
  \omega_{hk}=\exp\!\left[2\pi i(hs_x+ks_y)\right]
  =\exp\!\left[\frac{2\pi i}{3}(h+2k)\right],
  \label{eq:pbi2_registry_phase}
\end{equation}
while the two local orientations carry form factors $F_+(h,k,L)$ and
$F_-(h,k,L)$.

Adjacent trilayers obey the first-order interface law
\begin{equation}
  \boldsymbol{\theta}=(a,b_+,b_-,d_+,d_-),\qquad
  a+b_++b_-+d_++d_-=1.
  \label{eq:pbi2_transition_probabilities}
\end{equation}
Here $a$ denotes no registry or orientation change. $b_+$ and $b_-$ are
forward and backward registry shifts without an orientation change.
$d_+$ and $d_-$ are the two handed shift-plus-orientation-change events.  A
change of orientation without a registry shift is excluded.  The selected
deterministic limits are $a=1$ for 2H, $d_+=1$ for 4H, and $b_+=1$ for 6H.
The opposite registry directions generate their symmetry-related handed
counterparts.  Figure~\ref{fig:pbi2_transition_states} renders these cycles as
projected I--Pb--I stacks. The sequences are read from bottom to top: $0F_+\rightarrow0F_+$ for 2H,
$0F_+\rightarrow1F_-\rightarrow0F_+$ for 4H, and
$0F_+\rightarrow1F_+\rightarrow2F_+\rightarrow0F_+$ for 6H.  These parent
structures and PbI$_2$ polytypism are established experimentally.\cite{Minagawa1975,Palosz1990}

\begin{figure}[H]
	\centering
	\resizebox{0.88\linewidth}{!}{%
		\input{figures/results_pbi2/transition_matrix/pbi2_polytype_stacks.tex}%
	}
	\caption{Projected I--Pb--I trilayers for the deterministic PbI$_2$ parent
		cycles used in the transition model. (a) 2H repeats $0F_+$ under $a=1$.
		(b) An isolated $F_+$ trilayer shows the Pb and I atomic planes used to
		define the local orientation. (c) 4H alternates $0F_+$ and $1F_-$ under
		$d_+=1$. (d) 6H advances
		$0F_+\rightarrow1F_+\rightarrow2F_+\rightarrow0F_+$ under $b_+=1$.
		The $c$ axis points upward, so each stack is read from bottom to top. In a
		state label $rF_\sigma$, $r\in\{0,1,2\}$ is the basal registry and
		$\sigma\in\{+,-\}$ is the trilayer orientation. Violet and gray spheres
		denote I and Pb, green vertical guides mark lateral registry, blue
		horizontal guides mark atomic planes, and brackets group each I--Pb--I
		trilayer. Signs are shown on the reference $F_+$ trilayer.}
	\label{fig:pbi2_transition_states}
\end{figure}

The general layer-correlation treatment follows the established
transition-matrix theory of one-dimensionally disordered crystals, whereas the
six states and allowed PbI$_2$ interface events above are the material-specific
construction used here.\cite{HendricksTeller1942,KakinokiKomura1954,KakinokiKomura1965,Hart2018}
Because the registry labels enter a given rod only through
$1$, $\omega_{hk}$, and $\omega_{hk}^{-1}$, the full six-state matrix reduces
exactly to the following phase-weighted matrix in the orientation basis:
\begin{equation}
  M(\omega)=
  \begin{pmatrix}
    a+b_+\omega+b_-\omega^{-1} & d_+\omega+d_-\omega^{-1}\\
    d_+\omega^{-1}+d_-\omega & a+b_+\omega+b_-\omega^{-1}
  \end{pmatrix}.
  \label{eq:pbi2_reduced_transition}
\end{equation}
The probability matrix $P=M(1)$ propagates the orientation populations, while
$M(\omega)^\Delta$ sums the probabilities and registry phases of all paths
connecting two trilayers separated by $\Delta$ interfaces.  The complete construction and its  reduction are given in Supplementary
Note S3..

For each parent-rich population $j\in\{2H,4H,6H\}$, the disorder parameter has
a direct probabilistic meaning:
\begin{equation}
  p_j(\text{parent event})=1-\epsilon_j,
  \qquad
  p_j(\text{each of the four alternative events})=\frac{\epsilon_j}{4}.
  \label{eq:pbi2_fault_parameterization}
\end{equation}
For 2H, 4H, and 6H the parent event is $a$, $d_+$, and $b_+$,
respectively.  Written in the order
$(a,b_+,b_-,d_+,d_-)$, the five probabilities in
Eq.~\eqref{eq:pbi2_fault_parameterization} form
$\boldsymbol{\theta}_j$ and are inserted directly into $M_j(\omega)$ and
$P_j=M_j(1)$.  Thus $\epsilon_j$ changes the matrix powers that determine how much phase coherence remains
between increasingly distant trilayers.

For one definite sequence $s_n=(r_n,\sigma_n)$, the layer amplitude is
$g_{s_n}=\omega^{r_n}F_{\sigma_n}$ and the coherent stack amplitude is
$A_N(L)=\sum_{n=0}^{N-1}\xi(L)^n g_{s_n}$, where
$\xi(L)=\exp(iq_zd)$ is the vertical phase between neighboring trilayer
centers.  Expanding $\langle|A_N|^2\rangle$ leaves two familiar
contributions: the intensity scattered by each trilayer individually and the
interference between distinct trilayer pairs.

For population $j$, define the orientation population of trilayer $m$ as
\begin{equation}
  \boldsymbol{\pi}_{j,m}
  =\left(p_{j,m}^{+},p_{j,m}^{-}\right)
  =\boldsymbol{\pi}_{j,0}P_j^m,
  \qquad p_{j,m}^{+}+p_{j,m}^{-}=1.
  \label{eq:pbi2_orientation_population}
\end{equation}
Here $j$ identifies the 2H-, 4H-, or 6H-rich population and $m$ identifies a
trilayer within the finite stack.  The quantities $p_{j,m}^{+}$ and
$p_{j,m}^{-}$ are ordinary probabilities that trilayer $m$ has orientation
$F_+$ or $F_-$, respectively; they are neither scattering amplitudes nor the
sample weights $w_j$ used below.

The disorder-averaged interference amplitude for trilayers $m$ and
$m+\Delta$ is
\begin{equation}
  K_{m,\Delta}^{(j)}=
  \begin{pmatrix}p_{j,m}^{+} & p_{j,m}^{-}\end{pmatrix}
  \begin{pmatrix}F_+^* & 0\\0 & F_-^*\end{pmatrix}
  M_j(\omega)^\Delta
  \begin{pmatrix}F_+\\F_-\end{pmatrix}.
  \label{eq:pbi2_pair_kernel}
\end{equation}
The finite-stack intensity is then written as ``self scattering plus pair
interference'':
\begin{subequations}
\label{eq:pbi2_finite_intensity}
\begin{align}
  I_j(L)
  &=I_{j,\mathrm{self}}(L)+I_{j,\mathrm{pair}}(L),
  \label{eq:pbi2_intensity_split}\\
  I_{j,\mathrm{self}}(L)
  &=\frac{1}{N}\sum_{m=0}^{N-1}
  \left[p_{j,m}^{+}|F_+|^2+p_{j,m}^{-}|F_-|^2\right],
  \label{eq:pbi2_self_intensity}\\
  I_{j,\mathrm{pair}}(L)
  &=\frac{2}{N}\operatorname{Re}
  \sum_{\Delta=1}^{N-1}\xi(L)^\Delta
  \sum_{m=0}^{N-1-\Delta}K_{m,\Delta}^{(j)}.
  \label{eq:pbi2_pair_intensity}
\end{align}
\end{subequations}
The self term averages the isolated-trilayer intensity over the two possible
orientations.  The pair term groups all trilayer pairs by their separation
$\Delta$: there are $N-\Delta$ pairs at that separation, and each carries the
same vertical phase $\xi(L)^\Delta$.  In
Eq.~\eqref{eq:pbi2_pair_kernel}, the expression can be read from right to
left.  The final column supplies the scattering amplitude of the later
trilayer, $M_j(\omega)^\Delta$ sums all allowed stacking paths and registry
phases between the two layers, the diagonal matrix supplies the conjugated
amplitude of the earlier trilayer, and $\boldsymbol{\pi}_{j,m}$ averages over
whether that earlier trilayer is in the $F_+$ or $F_-$ orientation.  The
factor $2\operatorname{Re}$ in Eq.~\eqref{eq:pbi2_pair_intensity} includes the
reverse-ordered conjugate pair without writing it separately.  The complete
derivation is given in Supplementary Note S3.

The 2H-, 4H-, and 6H-rich domains are taken to be mutually incoherent, so the
three calculated intensities (not their amplitudes) are combined as
\begin{equation}
  I_{\mathrm{calc}}(L)=S\left\{\left[
  w_{2H}I_{2H}(L)+w_{4H}I_{4H}(L)+w_{6H}I_{6H}(L)
  \right]\!\otimes R(L)\right\},
  \qquad w_{2H}+w_{4H}+w_{6H}=1.
  \label{eq:pbi2_population_sum}
\end{equation}
Each $I_j(L)$ is first evaluated from
Eq.~\eqref{eq:pbi2_finite_intensity} using its own $\epsilon_j$.  The weights
then specify the fraction of the total calculated intensity assigned to each
parent-rich population, the shared resolution function $R(L)$ broadens the
combined profile, and $S$ converts the normalized calculation to the measured
intensity scale.  Accordingly, $\epsilon_j$ is an effective per-interface
disorder probability under the stated equal-alternative template, whereas
$w_j$ is an effective population weight.

For $h=k=0$, $\omega_{hk}=1$ and $F_+=F_-$.  Registry shifts and orientation
changes then leave the scattering amplitude unchanged, so stacking faults are
contrast-matched rather than absent.  The oscillations on the $m=0$ rod are
therefore finite-thickness fringes.  The same criterion applies to any other
rod for which both the registry phase and the orientation form factors are
indistinguishable.

\subsection{Refinement}

The disorder stage is applied only after the detector geometry, sample
alignment, mosaic distribution, lattice constants, ordered PbI$_2$ structure
factor, and profile resolution have been fixed by the upstream refinement.
Measured and calculated detector images are projected through the same
constant-$Q_R$ rod masks and reported versus $Q_z$ or the corresponding
stacking coordinate $L$.  No independent peak areas are introduced before the
stacking model is evaluated, because doing so would absorb the diffuse
between-peak intensity that the transition law is intended to explain.

The core disorder parameters are
\begin{equation}
  \boldsymbol{\Theta}_{\mathrm{dis}}=
  (\epsilon_{2H},\epsilon_{4H},\epsilon_{6H},
  w_{2H},w_{4H},w_{6H},N),
  \qquad w_{2H}+w_{4H}+w_{6H}=1,
  \label{eq:pbi2_disorder_parameter_vector}
\end{equation}
so the three weights contain only two independent degrees of freedom.  A
stable refinement can first impose a shared disorder magnitude
$\epsilon_{2H}=\epsilon_{4H}=\epsilon_{6H}$ and release the three values only
when the retained fault-sensitive rods support phase-dependent broadening.
The selected-profile residual is
\begin{equation}
  r_j=\frac{I_j^{\mathrm{meas}}-
  I_j^{\mathrm{model}}(\boldsymbol{\Theta}_{\mathrm{dis}})}{\sigma_j},
  \label{eq:pbi2_disorder_profile_residual}
\end{equation}
with any allowed global scale or low-order background terms declared
separately.  The fit is accepted only when it improves the selected diffuse
profiles without degrading the ordered detector-space agreement established
in the preceding stages.

\begin{figure}[H]
  \centering
  {\large\bfseries Disorder Series Fitting\par}
  \vspace{0.8em}
  \begin{minipage}[t]{0.32\linewidth}
    \centering
    \paneltagged{(a)}{\includegraphics[width=\linewidth]{figures/results_pbi2/low_disorder/Y2_5d_10m.png}}

    \vspace{0.6em}

    \paneltagged{(b)}{\includegraphics[width=\linewidth]{figures/results_pbi2/low_disorder/figure7_pbi2_qr_rod_qz_profiles.png}}
  \end{minipage}
  \hfill
  \begin{minipage}[t]{0.32\linewidth}
    \centering
    \paneltagged{(c)}{\includegraphics[width=\linewidth]{figures/results_pbi2/med_disorder/mid_disorder_detector_selected_q_regions_4deg.png}}

    \vspace{0.6em}

    \paneltagged{(d)}{\includegraphics[width=\linewidth]{figures/results_pbi2/med_disorder/figure7_pbi2_qr_rod_qz_profiles.png}}
  \end{minipage}
  \hfill
  \begin{minipage}[t]{0.32\linewidth}
    \centering
    \paneltagged{(e)}{\includegraphics[width=\linewidth]{figures/results_pbi2/high_disorder/Clean2_4d_2m.png}}

    \vspace{0.6em}

    \paneltagged{(f)}{\includegraphics[width=\linewidth]{figures/results_pbi2/high_disorder/figure7_pbi2_qr_rod_qz_profiles.png}}
  \end{minipage}

  \caption[PbI$_2$ diffuse-rod validation across disorder]{
  PbI$_2$ diffuse-rod comparison at $\theta_i=4^\circ$ for samples spanning
  increasing stacking disorder.  The top row shows measured detector images
  and the constant-$Q_R$ extraction regions.  Solid curves mark the projected
  rod centerlines, shaded bands mark the integrated detector pixels, and
  dashed outlines mark the integration boundaries.  The bottom row compares
  profiles obtained by applying the same projection to data and calculation;
  black solid curves are measured and orange dashed curves are calculated.
  The series tests whether one detector-space geometry and ordered baseline,
  augmented by the transition-matrix stacking model, reproduce the progressive
  transfer of intensity from ordered maxima into the diffuse rod signal.}
  \label{fig:pbi2_diffuse_rod_validation_disorder_series}
\end{figure}

Figure~\ref{fig:pbi2_diffuse_rod_validation_disorder_series} shows the visible
progression from sharp ordered maxima to stronger between-peak intensity.  The
absence of comparable diffuse contrast on $m=0$ does not imply that those
scattering volumes contain no faults; it follows from the phase-insensitive
condition described above.  Likewise, changes in the visibility of higher
order $m=0$ peaks are assigned to the growth-dependent mosaic and resolution
state rather than to lateral-registry disorder.

The ordered baseline and the transition-model parameterization are summarized
in Table~\ref{tab:pbi2_sf_parameters}.  Numerical values from the superseded
scalar two-population fit are not carried into the new model because they do
not map uniquely onto three population-specific $\epsilon_j$ values and three
normalized weights.

\begin{table}[H]
  \centering
  \caption{PbI$_2$ ordered-baseline parameters and stacking-disorder
  parameterization for the selected-rod comparison.  Numerical
  transition-matrix parameters must be reported from a refit using
  Eqs.~\eqref{eq:pbi2_fault_parameterization}--\eqref{eq:pbi2_population_sum};
  values from the superseded scalar model are not reused.}
  \label{tab:pbi2_sf_parameters}
  \scriptsize
  \setlength{\tabcolsep}{3pt}
  \renewcommand{\arraystretch}{1.12}
  \begin{tabularx}{\linewidth}{@{}>{\raggedright\arraybackslash}p{2.1cm} >{\raggedright\arraybackslash}X >{\raggedright\arraybackslash}X >{\raggedright\arraybackslash}X@{}}
    \toprule
    Parameter class & Fitted value or model quantity & Literature comparison & Role in the selected-rod model \\
    \midrule
    Lattice and coordinates &
    $a=4.5583\,\text{\AA}$, $c=6.985\,\text{\AA}$; Pb at $(0,0,0)$; I1 at $(1/3,2/3,0.2677)$ and I2 at $(2/3,1/3,0.7323)$. &
    Lattice constants remain within the 2H single-crystal structural spread, and the iodine coordinate follows the corrected $z\simeq0.268$ value.\cite{Minagawa1975,Palosz1990} &
    Fixes the ordered Bragg positions and the $|F_{hk}(L)|^2$ envelope before stacking disorder is introduced. \\
    Occupancies &
    A weakly floating test returned $o_{\mathrm{Pb}}=0.999$, $o_{\mathrm{I1}}=1.001$, and $o_{\mathrm{I2}}=1.000$; the ordered baseline was then evaluated with full occupancies. &
    The clean ordered model uses full occupancies; the density-deficient occupancy reported for a Palosz-type sample is sample-specific and was not used as the default model.\cite{Palosz1990,Lin2019} &
    Keeps the PbI$_2$ average structure stoichiometric while the diffuse rod intensity is assigned to stacking correlations. \\
    Directional ADPs &
    $U_R/U_z=0.0105/0.0185$ for Pb and $0.016/0.024\,\text{\AA}^2$ for I1 and I2. &
    Site-resolved Pb and I ADPs are not widely or consistently reported for direct transfer; the qualitative inequality $U_z(\mathrm{Pb})>U_R(\mathrm{Pb})$ is retained.\cite{Minagawa1975,Trueblood1996} &
    Moderates relative ordered intensities and is interpreted as a sample-specific displacement refinement rather than a universal PbI$_2$ displacement constant. \\
    Transition-matrix stacking disorder &
    $\epsilon_{2H}$, $\epsilon_{4H}$, $\epsilon_{6H}$; normalized weights $w_{2H}$, $w_{4H}$, $w_{6H}$; and finite stack length $N$.  Final numerical values are pending the transition-matrix refit. &
    These are effective selected-rod descriptors under the stated parent rules and equal-alternative fault templates, not universal PbI$_2$ constants. &
    Controls the redistribution of the ordered rod envelope into broadened maxima, finite-thickness fringes, and diffuse between-peak intensity. \\
    \bottomrule
  \end{tabularx}
\end{table}

The physical interpretation is therefore deliberately conditional on the
model.  Each $\epsilon_j$ measures the probability of departing from the
selected parent event within population $j$, the $w_j$ describe the incoherent
population mixture, and $N$ controls the finite-stack interference.  Numerical
values should be inserted only after the experimental profiles have been
refit with this parameterization.
